\chapter{Strabismus Surgery}

\section*{Introduction \& Clinical Indications}
Strabismus surgery is a specialized ophthalmic procedure aimed at correcting the misalignment of the eyes, a condition commonly referred to as strabismus. This misalignment can manifest as esotropia, where one or both eyes turn inward, or exotropia, where the eyes turn outward. The surgery involves the adjustment of the extraocular muscles responsible for eye movement, allowing for the repositioning, weakening, or strengthening of these muscles to achieve proper alignment. The clinical significance of strabismus surgery extends beyond cosmetic improvement; it plays a crucial role in restoring binocular vision, enhancing depth perception, and preventing amblyopia, particularly in pediatric patients. The procedure is applicable to both children and adults, with the overarching goal of achieving coordinated eye movement and improving the overall quality of life for patients suffering from this condition.

\begin{itemize}
  \item Eye muscle surgery.
  \item Extraocular muscle surgery.
  \item Strabismus correction.
  \item Squint surgery.
  \item Recession and resection procedures.
  \item Adjustable suture strabismus surgery.
\end{itemize}

The surgical principle underlying strabismus surgery is to achieve a balance in the forces exerted by the extraocular muscles, thereby aligning the eyes in primary gaze and ensuring coordinated movement. The operative techniques involve either weakening an overacting muscle through recession, where the muscle's insertion is moved posteriorly, or strengthening an underacting muscle through resection, where a portion of the muscle is excised and the remaining muscle is reattached to the globe. 

In cases of horizontal strabismus, adjustments are made to the medial and lateral rectus muscles to correct esotropia or exotropia, while vertical strabismus involves the vertical recti and oblique muscles. The ultimate goals of the surgery are to achieve long-lasting alignment, preserve full range of motion, and enhance patient comfort.

Strabismus surgery has a rich history, dating back to the mid-nineteenth century when early procedures involved tenotomy of the medial rectus muscle for the treatment of esotropia. Pioneers such as Sir Harold Gillies and Dr. William Wright contributed significantly to the development of surgical techniques in the early 1900s. 

By the mid-twentieth century, advancements in the understanding of ocular motility led to the establishment of standardized recession and resection techniques. The introduction of adjustable suture techniques in the 1970s marked a significant evolution in strabismus surgery, allowing for intraoperative adjustments to fine-tune muscle positioning. Today, modern strabismus surgery benefits from enhanced anesthesia protocols, precise preoperative measurements, and refined surgical techniques, resulting in improved outcomes for patients.

Strabismus surgery is indicated in various clinical scenarios, including:

\begin{itemize}
  \item Persistent esotropia or exotropia not corrected by glasses, patching, or vision therapy.
  \item Strabismus causing or risking amblyopia, or lazy eye, in children.
  \item Loss of binocular vision and depth perception due to misalignment.
  \item Double vision (diplopia) in adults from acquired strabismus.
  \item Misalignment appearing after trauma, stroke, or thyroid eye disease.
  \item Cosmetic and functional correction of a visible eye turn.
  \item Nystagmus-related or paralytic strabismus when other treatment fails.
  \item Significant deviation impacting quality of life or social interactions.
\end{itemize}

Standardized diagnostic scores, such as the Alvarado Score for appendicitis, are not typically applied in strabismus but clinical assessment remains critical.

Strabismus surgery is primarily a corrective procedure for structural misalignment that has not resolved with conservative measures. It is often considered a secondary intervention after optical correction, patching, and vision therapy have been attempted, particularly in children with accommodative esotropia. In adults with acquired strabismus due to conditions such as thyroid eye disease, surgery is usually performed once the underlying disease is stable. The timing and choice of surgical technique are tailored to the specific type and magnitude of the deviation, the patient's age, and the potential for residual binocular vision.

\section*{Relevant Surgical Anatomy}
The anatomy relevant to strabismus surgery includes the extraocular muscles, which consist of six muscles that control eye movement: the medial and lateral rectus, superior and inferior rectus, and superior and inferior oblique muscles. These muscles are innervated primarily by the oculomotor nerve (cranial nerve III), with the lateral rectus being innervated by the abducens nerve (cranial nerve VI) and the superior oblique by the trochlear nerve (cranial nerve IV). 

The arterial supply to these muscles is primarily from the ophthalmic artery, a branch of the internal carotid artery, while venous drainage occurs through the superior and inferior ophthalmic veins. Important surgical landmarks include:

\begin{itemize}
  \item McBurney's point for lateral rectus muscle access.
  \item Calot's triangle for identifying the anatomical relationships of the extraocular muscles.
  \item The tendons of the extraocular muscles, which are critical for muscle recession and resection techniques.
\end{itemize}

Understanding these anatomical relationships is essential for minimizing complications and ensuring successful surgical outcomes.

\section*{Preoperative Workup \& Preparation}
Prior to surgery, a comprehensive preoperative evaluation is essential. This includes a complete eye examination with careful measurement of the angle of deviation in various gaze positions, assessment of binocular function, and review of any prior treatments such as patching or glasses. 

The patient is typically advised to fast prior to general anesthesia. Adjustments to anticoagulant and antiplatelet medications may be necessary. The surgeon discusses the risks of under- or over-correction, which may necessitate further surgical intervention. Patients should also arrange for postoperative assistance and be advised to avoid rubbing the eye during recovery.

The contraindications and precautions for strabismus surgery can be categorized as follows:

Absolute contraindications:

\begin{itemize}
  \item Uncontrolled amblyopia that limits functional outcomes.
  \item Active thyroid eye disease or other conditions that may affect muscle stability.
  \item Severe systemic illness that poses a risk during anesthesia.
\end{itemize}

Relative contraindications:

\begin{itemize}
  \item Recent ocular trauma or surgery that may complicate healing.
  \item Patients with limited eye movement, requiring tailored surgical planning.
  \item Unstable refractive errors that may affect alignment postoperatively.
\end{itemize}

Precautions include careful estimation of the amount of recession or resection to avoid significant over- or under-correction, as well as avoiding injury to the sclera and surrounding structures during dissection.

\section*{Surgical Technique \& Procedure}
Strabismus surgery is typically performed under general anesthesia for children, while adults may undergo the procedure under local or topical anesthesia if they are cooperative. The surgical steps include:

\begin{itemize}
  \item Administration of general or local anesthesia, with appropriate draping of the surgical field.
  \item Creation of a conjunctival incision to expose the target extraocular muscle.
  \item Isolation of the muscle using a muscle hook.
  \item Performing recession or resection of the appropriate muscles based on the preoperative plan.
  \item Reattachment of the muscle with sutures, ensuring proper tension.
  \item Optional adjustment of sutures if adjustable techniques are employed.
  \item Closure of the conjunctiva with sutures.
\end{itemize}

The procedure generally lasts one to two hours and may involve one or both eyes, depending on the extent of the strabismus.

\section*{Complications \& Risks}
The potential complications and risks associated with strabismus surgery can be categorized as follows:

General surgical complications:

\begin{itemize}
  \item Bleeding and infection at the surgical site.
  \item Anesthesia-related complications.
\end{itemize}

Procedure-specific risks:

\begin{itemize}
  \item Under-correction or over-correction of eye alignment.
  \item Residual or recurrent strabismus requiring additional surgery.
  \item Double vision, which may be transient or persistent.
  \item Decreased or altered eye movement postoperatively.
  \item Postoperative redness, swelling, and discomfort.
  \item Scleral perforation, a rare but serious complication.
  \item Anterior segment ischemia in extensive muscle surgery.
  \item Suture-related complications and conjunctival scarring.
\end{itemize}

Awareness of these risks is crucial for informed consent and patient management.

\section*{Postoperative Care \& Recovery}
Postoperative recovery following strabismus surgery typically involves several phases. In the immediate postoperative period, patients are monitored in the Post Anesthesia Care Unit (PACU) for stabilization. 

During the first 24-48 hours, patients may experience redness, swelling, and mild discomfort, which can be managed with prescribed analgesics. The use of artificial tears and antibiotic or steroid drops is common to promote healing and prevent infection. 

In the first week, patients are advised to avoid strenuous activities and rubbing the eyes. Most individuals can return to normal activities within a week, although full stabilization of eye alignment may take several weeks as swelling subsides and the muscles settle into their new positions.

During strabismus surgery, the surgeon documents the alignment of the eyes and the condition of the extraocular muscles. The pathology report on any resected muscle tissue typically indicates the presence of normal muscle fibers without significant pathological changes. The surgeon assesses the effectiveness of the procedure based on the alignment achieved and any residual deviations that may require further intervention.

A successful postoperative course is characterized by the resolution of pain, improvement in eye alignment, and the restoration of binocular function when applicable. Patients should expect gradual improvement in their symptoms, with alignment stabilizing over the weeks following surgery. 

Monitoring for any signs of complications, such as persistent double vision or significant misalignment, is essential during follow-up visits.

Postoperative care includes regular follow-up visits to assess eye alignment, healing, and overall visual function. 

Key aspects of postoperative care include:

\begin{itemize}
  \item Suture or staple removal as indicated.
  \item Assessment of vision and binocular function.
  \item Continued patching or glasses for children as needed.
  \item Adjustment of adjustable sutures in the immediate postoperative period.
  \item Monitoring for amblyopia and visual development in pediatric patients.
\end{itemize}

Patients should be educated on warning signs that necessitate contacting their healthcare provider, such as increased pain, significant swelling, or changes in vision.

\section*{Outcomes \& Prognosis}
Strabismus affects approximately 2 to 4 percent of the general population, making it one of the most prevalent pediatric ophthalmic conditions. Esotropia is more commonly observed in young children, often emerging in the first few years of life, while exotropia tends to become more apparent later in childhood. 

The condition has a genetic predisposition and is associated with refractive errors, particularly hyperopia, as well as prematurity, neurological disorders, and various syndromes. If left untreated, childhood strabismus can lead to amblyopia and permanent loss of binocular vision, underscoring the importance of early detection and intervention. Strabismus surgery remains one of the most frequently performed procedures in pediatric ophthalmology.

Strabismus surgery boasts a high success rate, with studies reporting satisfactory alignment in over 75 to 80 percent of patients following a single procedure. Success rates improve with the use of adjustable sutures and careful patient selection. 

Functional outcomes are influenced by factors such as the patient's age, the presence of binocular vision prior to surgery, and the underlying cause of the strabismus. Children who undergo early surgical intervention have the best chance of developing normal binocular vision and avoiding amblyopia, while adults with acquired strabismus often experience relief from double vision and improved cosmetic appearance. 

Some patients may require additional surgeries over time, emphasizing the need for ongoing follow-up and management.

\section*{References}
\begin{enumerate}
  \item MedlinePlus. Strabismus Surgery. U.S. National Library of Medicine; 2025.
  \item American Academy of Ophthalmology. Pediatric Ophthalmology and Strabismus. AAO; 2024.
  \item Current Medical Diagnosis and Treatment. McGraw-Hill; 2025.
  \item von Noorden GK, Campos EC. Binocular Vision and Ocular Motility. Mosby; 2023.
\end{enumerate}

\clearpage
